\section{Extended Experimental Protocol}

This appendix reports the full experimental protocol for CoRAS. All experiments use episode-block splits with disjoint tuning, calibration, and test partitions. The tuning split is used only for prompt/adaptor tuning and temperature scaling. The calibration split is used only to compute conformal thresholds. The test split is held out until final reporting.

\subsection{Datasets and Splits}

For each dataset we report the number of frames, episodes, action-token classes, domains, and train/tune/calibration/test counts. Continuous robot actions are quantized into action-token labels using KMeans; the codebook size and quantization MSE are reported for each public dataset.

\subsection{Methods}

We compare uncalibrated top-1, calibrated top-$k$, vanilla conformal prediction, APS conformal prediction, temperature-scaled conformal prediction, prompt-only top-1, CoRAS inverse-probability conformal prediction, CoRAS APS, and Mondrian/domain-conditional CoRAS.

\subsection{Main Aggregate Table}

The table below is generated by \texttt{scripts/aggregate\_results.py}. Replace the input path with the relevant track's aggregate table.

\begin{tabular}{llllll}
\toprule
Method & Coverage & |Coverage gap| & Mean set size & Top-1 acc. & Unsafe singleton error \\
\midrule
aps\_conformal & 0.918±0.026 & 0.026±0.016 & 9.863±1.274 & 0.471±0.082 & 0.000±0.001 \\
coras & 0.900±0.019 & 0.013±0.012 & 5.986±1.211 & 0.624±0.039 & 0.009±0.008 \\
coras\_aps & 0.942±0.011 & 0.042±0.011 & 9.068±0.434 & 0.624±0.039 & 0.000±0.000 \\
mondrian\_domain & 0.898±0.022 & 0.019±0.004 & 6.042±1.171 & 0.624±0.039 & 0.009±0.006 \\
prompt\_only\_top1 & 0.624±0.039 & 0.276±0.039 & 1.000±0.000 & 0.624±0.039 & 0.376±0.039 \\
temperature\_conformal & 0.899±0.027 & 0.017±0.019 & 8.931±2.485 & 0.471±0.082 & 0.004±0.008 \\
top1\_singleton & 0.471±0.082 & 0.429±0.082 & 1.000±0.000 & 0.471±0.082 & 0.529±0.082 \\
topk\_calibrated & 0.908±0.032 & 0.025±0.019 & 10.800±2.387 & 0.471±0.082 & 0.000±0.000 \\
vanilla\_conformal & 0.897±0.028 & 0.019±0.018 & 8.501±2.265 & 0.471±0.082 & 0.008±0.012 \\
\bottomrule
\end{tabular}


\subsection{Per-Domain Coverage}

Report \texttt{aggregate\_domain\_metrics.csv} for each track. The key reviewer-facing diagnostic is whether target-domain coverage remains near $1-\alpha$ and whether set-size inflation is concentrated in specific shifts.

\subsection{Calibration-Size Sweep}

Report \texttt{aggregate\_metrics.csv} grouped by method, alpha, and calibration fraction. This demonstrates whether CoRAS remains stable when calibration labels are scarce.

\subsection{Failure Cases}

Use \texttt{prediction\_sets\_coras\_alpha0.10.csv} to identify false singleton executions and large-set examples. Include a small figure grid for the most representative failure cases.
